\documentclass[12pt, a4paper, oneside, french]{report}

% Encodage et langue
\usepackage[utf-8]{inputenc}
\usepackage[T1]{fontenc}
\usepackage[french]{babel}

% Mathématiques et symboles
\usepackage{amsmath}
\usepackage{amssymb}
\usepackage{textcomp}

% Listings pour le code source
\usepackage{listings}
\usepackage{xcolor}
\usepackage{fancyvrb}

% Définition du style pythonstyle pour les listings
\lstdefinestyle{pythonstyle}{
    language=Python,
    basicstyle=\ttfamily\small,
    keywordstyle=\color{blue},
    commentstyle=\color{gray},
    stringstyle=\color{red},
    breaklines=true,
    breakatwhitespace=true,
    tabsize=4,
    showstringspaces=false,
    captionpos=b,
    frame=single,
    backgroundcolor=\color{white}
}

% Autres packages
\usepackage{graphicx}
\usepackage{booktabs}
\usepackage{hyperref}
\usepackage{geometry}

% Géométrie des pages
\geometry{
    left=2.5cm,
    right=2.5cm,
    top=2.5cm,
    bottom=2.5cm
}

% Titre du document
\title{Annexe — Corrections méthodologiques et contrôle des fuites de données}
\author{}
\date{}

\begin{document}

\maketitle

% ==============================================================================
% ANNEXE — Corrections méthodologiques des pipelines
% SHI, BiLSTM--CatBoost, SPI et détection des faux démarrages
% ==============================================================================

\chapter*{Annexe — Corrections méthodologiques et contrôle des fuites de données}
\addcontentsline{toc}{chapter}{Annexe — Corrections méthodologiques et contrôle des fuites de données}

Cette annexe présente les corrections méthodologiques apportées à plusieurs
pipelines expérimentaux développés dans le cadre de cette thèse. L'objectif
principal est de documenter les procédures mises en œuvre pour identifier,
corriger et contrôler les risques de fuite de données (\textit{data leakage})
susceptibles de surestimer les performances des modèles d'apprentissage
automatique.

Les trois pipelines principaux concernés sont :

\begin{itemize}
    \item la construction de l'indice de santé des sols (SHI) à partir des
    variables biologiques du sol ;
    \item la prédiction des rendements par une architecture
    BiLSTM--CatBoost et l'évaluation comparative par le test de
    Diebold--Mariano ;
    \item la prédiction de l'indice standardisé des précipitations (SPI)
    à partir de variables satellitaires et climatiques.
\end{itemize}

Un contrôle complémentaire des données est également fourni afin de vérifier
automatiquement l'absence de variables susceptibles d'introduire une
information provenant directement de la variable cible dans les variables
explicatives.

% ==============================================================================
\section*{A. Pipeline SHI : correction du risque de fuite de données}
\addcontentsline{toc}{section}{A. Pipeline SHI : correction du risque de fuite de données}

\subsection*{A.1. Construction de l'indice SHI}

L'indice de santé des sols (\textit{Soil Health Index}, SHI) est construit
par agrégation pondérée de six indicateurs biologiques :

\begin{equation}
SHI =
0.4\,Cmic
+0.1\,BR
+0.1\,Xylosidase
+0.1\,\beta\text{-glucosidase}
+0.1\,Acid\_phosphatase
+0.2\,N\text{-acetylglucosaminidase}.
\end{equation}

où :

\begin{itemize}
    \item $Cmic$ représente le carbone de la biomasse microbienne ;
    \item $BR$ représente la respiration basale ;
    \item $Xylosidase$ représente l'activité de la xylosidase ;
    \item $\beta$-glucosidase représente l'activité de la
    $\beta$-glucosidase ;
    \item $Acid\_phosphatase$ représente l'activité de la phosphatase acide ;
    \item $N$-acetylglucosaminidase représente l'activité de la
    $N$-acétylglucosaminidase.
\end{itemize}

Les poids utilisés sont issus de la procédure de sélection décrite dans le
manuscrit et permettent d'obtenir un indice composite représentant l'état
biologique du sol.

\subsection*{A.2. Correction du protocole expérimental}

La procédure corrigée impose que toutes les transformations susceptibles
d'utiliser des informations statistiques soient réalisées à partir des
données d'entraînement uniquement.

Le jeu de données corrigé contient :

\begin{itemize}
    \item 881 observations au total ;
    \item 704 observations pour l'entraînement ;
    \item 177 observations pour le test ;
    \item une proportion de 80\,\% pour l'entraînement et 20\,\% pour le test ;
    \item \texttt{random\_state = 42}.
\end{itemize}

Le modèle final retenu après comparaison des algorithmes est
\textbf{CatBoost}.

Les performances obtenues sur le jeu de test sont :

\begin{table}[htbp]
\centering
\caption{Performances du modèle CatBoost pour la prédiction du SHI}
\begin{tabular}{lc}
\hline
Métrique & Valeur \\
\hline
$R^2$ test & 0.850 \\
MAE & 51.2 \\
RMSE & 80.0 \\
$R^2$ validation croisée & 0.819 \\
\hline
\end{tabular}
\end{table}

Le contrôle final du pipeline indique que les variables explicatives ne
contiennent pas directement la variable cible ou une transformation de
celle-ci susceptible de provoquer une fuite de données.

\subsection*{A.3. Code source corrigé}

Le code Python correspondant est fourni dans le fichier :

\begin{center}
\texttt{shi\_lucas2018\_leakage\_free.py}
\end{center}

\lstinputlisting[
    style=pythonstyle,
    caption={Pipeline SHI corrigé et sans fuite de données},
    label={lst:shi-leakage-free}
]{shi_lucas2018_leakage_free.py}


% ==============================================================================
\section*{B. Pipeline BiLSTM--CatBoost et test de Diebold--Mariano}
\addcontentsline{toc}{section}{B. Pipeline BiLSTM--CatBoost et test de Diebold--Mariano}

\subsection*{B.1. Objectif de la correction}

Une seconde vérification méthodologique a été réalisée sur le pipeline
associant un réseau BiLSTM à un modèle CatBoost.

L'objectif est de garantir que les variables utilisées pendant l'entraînement
ne contiennent aucune information provenant de la période ou de la valeur
cible à prédire.

La séparation temporelle des observations est conservée afin d'éviter qu'une
information provenant du futur soit utilisée pour prédire une observation
antérieure.

\subsection*{B.2. Comparaison des modèles}

Après correction du protocole, les prédictions obtenues par les différents
modèles peuvent être comparées à partir de leurs erreurs individuelles.

Le test de Diebold--Mariano est utilisé pour déterminer si la différence de
performance prédictive entre deux modèles est statistiquement significative.

Pour deux modèles $A$ et $B$, on définit la différence des pertes :

\begin{equation}
d_t = L(e_{A,t}) - L(e_{B,t}),
\end{equation}

où $L$ représente la fonction de perte et $e_{A,t}$ et $e_{B,t}$ les erreurs
de prédiction des deux modèles à la période $t$.

L'hypothèse nulle du test est :

\begin{equation}
H_0 : E[d_t] = 0.
\end{equation}

Elle signifie qu'il n'existe pas de différence significative entre les
performances prédictives des deux modèles considérés.

\subsection*{B.3. Code source corrigé}

Le code correspondant au pipeline BiLSTM--CatBoost et au test de
Diebold--Mariano est fourni dans :

\begin{center}
\texttt{bilstm\_catboost\_dm\_corrected\_utf8.py}
\end{center}

\lstinputlisting[
    style=pythonstyle,
    caption={Pipeline BiLSTM--CatBoost corrigé et test de Diebold--Mariano},
    label={lst:bilstm-catboost-dm}
]{bilstm_catboost_dm_corrected_utf8.py}


% ==============================================================================
\section*{C. Pipeline SPI : correction de la fuite de données}
\addcontentsline{toc}{section}{C. Pipeline SPI : correction de la fuite de données}

\subsection*{C.1. Problématique}

Le pipeline de prédiction du SPI a également fait l'objet d'un audit
méthodologique.

Une attention particulière a été portée à la construction des variables
explicatives, notamment aux variables climatiques et satellitaires calculées
sur des périodes susceptibles de recouvrir la période de calcul de la cible.

L'objectif de la correction est de garantir que les variables utilisées pour
la prédiction du SPI sont disponibles au moment où la prédiction est censée
être réalisée.

\subsection*{C.2. Formulation du SPI}

Dans le cadre du jeu de données considéré, la formulation utilisée après
standardisation peut être exprimée sous la forme :

\begin{equation}
SPI = \frac{X-\mu}{\sigma},
\end{equation}

où $X$ représente la quantité de précipitations considérée, $\mu$ sa moyenne
calculée sur la période de référence et $\sigma$ son écart-type.

\subsection*{C.3. Performances après correction}

Les performances obtenues après correction sont présentées dans le tableau
suivant.

\begin{table}[htbp]
\centering
\caption{Performances des modèles pour la prédiction du SPI après correction}
\begin{tabular}{lrrrrrr}
\hline
Modèle &
MSE$_{train}$ &
RMSE$_{train}$ &
$R^2_{train}$ &
MSE$_{test}$ &
RMSE$_{test}$ &
$R^2_{test}$ \\
\hline
RF &
0.04380 &
0.20927 &
0.9583 &
0.13866 &
0.37238 &
0.8205 \\
XGBoost &
0.01429 &
0.11956 &
0.9864 &
0.18860 &
0.43429 &
0.7559 \\
SVR &
0.18426 &
0.42925 &
0.8246 &
0.32597 &
0.57094 &
0.5781 \\
\hline
\end{tabular}
\end{table}

Les écarts entre les coefficients de détermination d'entraînement et de test
sont respectivement :

\begin{table}[htbp]
\centering
\caption{Écart entre les performances d'entraînement et de test}
\begin{tabular}{lc}
\hline
Modèle & $R^2_{train}-R^2_{test}$ \\
\hline
RF & 0.1378 \\
XGBoost & 0.2305 \\
SVR & 0.2465 \\
\hline
\end{tabular}
\end{table}

Le modèle Random Forest présente ainsi le meilleur compromis entre
performance prédictive et écart entre les performances d'entraînement et de
test parmi les modèles évalués.

\subsection*{C.4. Code source corrigé}

Le fichier actuellement présent dans le dossier \texttt{code} est :

\begin{center}
\texttt{spi\_prediction\_corrige (3).py}
\end{center}

Pour faciliter la compilation LaTeX et la gestion du dépôt GitHub, il est
fortement recommandé de renommer ce fichier en :

\begin{center}
\texttt{spi\_prediction\_corrige.py}
\end{center}

Après renommage, utiliser :

\lstinputlisting[
    style=pythonstyle,
    caption={Pipeline SPI corrigé et sans fuite de données},
    label={lst:spi-leakage-free}
]{spi_prediction_corrige.py}


% ==============================================================================
\section*{D. Détection des faux démarrages}
\addcontentsline{toc}{section}{D. Détection des faux démarrages}

Une procédure complémentaire concerne la détection des faux démarrages
phénologiques à partir des données pluviométriques et agroclimatiques.

Le pipeline corrigé est disponible dans le fichier :

\begin{center}
\texttt{version\_false\_start\_corrige\_sans\_leakage.py}
\end{center}

\lstinputlisting[
    style=pythonstyle,
    caption={Pipeline corrigé de détection des faux démarrages sans fuite de données},
    label={lst:false-start}
]{version_false_start_corrige_sans_leakage.py}


% ==============================================================================
\section*{E. Contrôle automatique des fuites de données}
\addcontentsline{toc}{section}{E. Contrôle automatique des fuites de données}

Un script indépendant a également été développé afin de contrôler les
variables utilisées dans les différents pipelines.

Le fichier est :

\begin{center}
\texttt{Controler\_data\_leakage.py}
\end{center}

Il permet notamment de rechercher les variables susceptibles d'être
directement ou indirectement liées à la variable cible et de vérifier la
cohérence entre les variables explicatives et la cible.

\lstinputlisting[
    style=pythonstyle,
    caption={Script de contrôle des risques de fuite de données},
    label={lst:control-leakage}
]{Controler_data_leakage.py}


% ==============================================================================
\section*{F. Organisation des codes sources}
\addcontentsline{toc}{section}{F. Organisation des codes sources}

Les scripts utilisés pour reproduire les expériences sont regroupés dans le
même répertoire que la présente annexe :

\begin{verbatim}
code/
|
|-- annexe_shi_spi_leakage_free.tex
|
|-- shi_lucas2018_leakage_free.py
|-- bilstm_catboost_dm_corrected_utf8.py
|-- spi_prediction_corrige.py
|-- version_false_start_corrige_sans_leakage.py
|-- Controler_data_leakage.py
\end{verbatim}

Cette organisation permet d'associer directement chaque correction
méthodologique à son implémentation informatique correspondante.

Les versions corrigées sont destinées à assurer une évaluation plus fiable
de la capacité de généralisation des modèles et à faciliter la reproductibilité
des résultats présentés dans la thèse.

\end{document}
